\documentclass[11pt]{article}
\usepackage[T1]{fontenc}
\usepackage{textcomp}
\usepackage[margin=0.75in]{geometry}
\usepackage{amsmath,amssymb,amsthm}
\usepackage{array,booktabs}
\usepackage{hyperref}
\setlength{\parindent}{0pt}
\newtheorem{theorem}{Theorem}
\theoremstyle{definition}
\newtheorem{definition}{Definition}
\title{Flow Proof Artifact: Nat-plus-mono-right}
\author{Generated by \texttt{flow doc proof}}
\date{Flow Proof Book}
\begin{document}
\maketitle
Adding the same term on the right preserves less-or-equal.\\[0.5em]
\textbf{Source.} peano — https://en.wikipedia.org/wiki/Monotonic\_function\\[0.5em]
\bigskip
\noindent{\large \textbf{Derived fact 1.} \textit{If b ≤ c then a + b ≤ a + c} \hfill the natural numbers \cdot addition \cdot adding on the right preserves order}\par
\smallskip\noindent\textit{Coordinate: the natural numbers \cdot addition \cdot adding on the right preserves order}\par
\smallskip\noindent\textit{Source: peano/induction — Gries \& Schneider, Ch. 3}\par
\smallskip\noindent\textit{Built on: less-or-equal is reflexive, for order on the natural numbers, less-or-equal is transitive, for order on the natural numbers, adding zero on the left does not change the number}\par
\medskip
\noindent\textbf{Goal.} If b ≤ c then a + b ≤ a + c\par
\noindent$\displaystyle \forall a \in \mathbb{N} \forall b \in \mathbb{N} \forall c \in \mathbb{N}\quad a + b \le a + c$\par
\medskip
\noindent
\begin{tabular}{@{} >{\raggedright\arraybackslash}p{0.06\textwidth} >{\raggedright\arraybackslash}p{0.47\textwidth} >{\raggedleft\arraybackslash}p{0.41\textwidth} @{}}
\textbf{\#} & \textbf{Proof} & \textbf{Mathematics} \\
\midrule
\textbf{1.} & We prove that adding on the right preserves order for addition on the natural numbers. & \\[0.45em]
\textbf{2.} & We proceed by induction on a: first the base case, then the inductive step. & \\[0.45em]
\textbf{3.} & Consider the base case in which b <= c. & \\[0.45em]
\textbf{4.} & Consider the base case in which a is zero. & \\[0.45em]
\textbf{5.} & We invoke the definitional clause governing addition on the natural numbers: adding zero on the left does not change the number (instantiated for b). & $\displaystyle 0 + b = b$ \\[0.45em]
\textbf{6.} & We invoke the definitional clause governing addition on the natural numbers: adding zero on the left does not change the number (instantiated for c). & $\displaystyle 0 + c = c$ \\[0.45em]
\textbf{7.} & From step 4, step 5, and step 6, we can deduce that a plus b is at most a plus c. This establishes the base case (see step 4, step 5, and step 6). Hence proven. & $\displaystyle a + b \le a + c$ \\[0.45em]
\textbf{8.} & For the inductive step, suppose the claim holds for all smaller values. & \\[0.45em]
\textbf{9.} & Under the supposition in step 8, let n denote the predecessor of a. & \\[0.45em]
\textbf{10.} & Under the supposition in step 8, we cross the inductive boundary: assume the claim holds for n (the induction hypothesis). & $\displaystyle n + b \le n + c$ \\[0.45em]
\textbf{11.} & We invoke the derived fact governing order on the natural numbers: less-or-equal is transitive, for order on the natural numbers (instantiated for n + b, c + b, c + c). & \\[0.45em]
\textbf{12.} & From step 8, step 9, step 10, and step 11, this implies a plus b is at most a plus c. Hence proven. & $\displaystyle a + b \le a + c$ \\[0.45em]
\bottomrule
\end{tabular}
\label{thm:Nat-addition-adding_on_the_right_preserves_order}
\par\medskip\hrule
\end{document}
